% Compile with XeLaTeX. Replace every \placeholder{...} before formal use.
\documentclass[UTF8,zihao=-4,fontset=fandol,a4paper]{ctexart}

\usepackage[top=25mm,bottom=25mm,left=25mm,right=25mm]{geometry}
\usepackage{amsmath,amssymb,bm}
\usepackage{booktabs,longtable,tabularx,array}
\usepackage{graphicx,float}
\usepackage{caption,subcaption}
\usepackage{xcolor}
\usepackage{listings}
\usepackage{fancyhdr}
\usepackage[hidelinks]{hyperref}
\usepackage[nameinlink,noabbrev]{cleveref}

\graphicspath{{figures/}{images/}}
\pagestyle{fancy}
\fancyhf{}
\cfoot{\thepage}
\renewcommand{\headrulewidth}{0pt}
\setlength{\parindent}{2em}
\setlength{\parskip}{0pt}
\linespread{1.35}

\ctexset{
  section = {
    name = {,、},
    number = \chinese{section},
    format = \centering\bfseries\zihao{3},
    beforeskip = 1.5ex,
    afterskip = 1.2ex
  },
  subsection = {
    format = \bfseries\zihao{4},
    beforeskip = 1.2ex,
    afterskip = 0.8ex
  },
  subsubsection = {
    format = \bfseries\zihao{-4},
    beforeskip = 1ex,
    afterskip = 0.6ex
  }
}

\captionsetup{font=small,labelsep=quad}
\lstset{
  basicstyle=\ttfamily\small,
  numbers=left,
  numberstyle=\tiny,
  frame=single,
  breaklines=true,
  showstringspaces=false,
  columns=fullflexible,
  keepspaces=true
}

\newenvironment{cumcmabstract}
  {\begin{center}\bfseries\zihao{4} 摘\quad 要\end{center}\noindent}
  {\par}
\newcommand{\keywords}[1]{\par\vspace{0.8ex}\noindent\textbf{关键词：}#1}
\newcommand{\placeholder}[1]{\textcolor{gray}{\textit{【#1】}}}

\title{\bfseries\zihao{3} \placeholder{论文题目}}
\date{}

\begin{document}

\maketitle
\thispagestyle{fancy}
\vspace{-3em}

\begin{cumcmabstract}
\placeholder{用一页以内说明问题、方法、最重要的定量结果、相对基线的改进及适用边界。摘要中的数字必须来自已经冻结的结果。}
\keywords{\placeholder{关键词1}；\placeholder{关键词2}；\placeholder{关键词3}}
\end{cumcmabstract}

\newpage

\section{问题重述与任务分解}

\placeholder{说明问题背景、每个子问题的交付物、变量关系和子问题之间的数据依赖。}

\section{问题分析与数据说明}

\placeholder{定义观测单位、字段含义、单位、缺失机制、异常处理依据和训练验证边界。}

\section{模型假设}

\begin{enumerate}
  \item \placeholder{假设及其现实依据；同时说明如何检查。}
  \item \placeholder{假设失效时对结论的影响。}
\end{enumerate}

\section{符号与单位}

\begin{table}[H]
  \centering
  \caption{主要符号与单位}
  \begin{tabularx}{0.88\textwidth}{>{\centering\arraybackslash}p{0.18\textwidth} X >{\centering\arraybackslash}p{0.18\textwidth}}
    \toprule
    符号 & 含义 & 单位 \\
    \midrule
    $x_i$ & \placeholder{第 $i$ 个决策变量的定义} & \placeholder{单位} \\
    \bottomrule
  \end{tabularx}
  \label{tab:symbols}
\end{table}

\section{模型建立与求解}

\subsection{基线模型}

\placeholder{先给出透明、可执行的基线，并说明其评价指标。}

\subsection{模型表述}

以一般优化问题为例，目标和约束可写为
\begin{equation}
  \min_{\bm{x}\in\mathcal{X}} f(\bm{x})
  \quad \text{s.t.} \quad
  g_j(\bm{x}) \le 0,\; j=1,\ldots,m .
  \label{eq:model}
\end{equation}
式\eqref{eq:model}中的每一项都应对应现实机制、数据字段或硬约束。

\subsection{求解过程与复现设置}

\placeholder{给出初始化、参数、停止条件、随机种子、软件版本、运行命令和计算复杂度。}

\section{结果与验证}

\subsection{主要结果与基线比较}

\placeholder{报告数据集或场景、指标、基线、数值和条件。不要只报告一次随机运行中的最好结果。}

\subsection{敏感性、稳健性与失败情形}

\placeholder{给出扰动范围、不确定性、可行性检查、失败区域以及结论边界。}

\section{结论与局限}

\placeholder{只总结有证据支持的贡献、决策含义和最重要的局限。}

\begin{thebibliography}{99}
  \bibitem{cumcm-format-2026}
  全国大学生数学建模竞赛组委会.
  全国大学生数学建模竞赛论文格式规范（2026年修订稿）[EB/OL].
  \url{https://www.mcm.edu.cn/html_cn/node/4cd596519c9eb9fbd866398f6df0caa3.html}.
\end{thebibliography}

\appendix
\ctexset{section={name={附录~,},number=\Alph{section}}}

\section{支撑材料文件清单}

\begin{longtable}{p{0.08\textwidth}p{0.35\textwidth}p{0.45\textwidth}}
  \toprule
  序号 & 文件路径 & 内容及复现作用 \\
  \midrule
  \endfirsthead
  \toprule
  序号 & 文件路径 & 内容及复现作用 \\
  \midrule
  \endhead
  1 & \placeholder{code/main.m 或 code/main.py} & \placeholder{主模型入口与运行命令} \\
  2 & \placeholder{results/canonical.mat} & \placeholder{论文采用的冻结结果及校验方式} \\
  \bottomrule
\end{longtable}

\section{完整源程序}

\placeholder{在此附上完整、可运行的源程序；较长程序可使用 \texttt{\textbackslash lstinputlisting} 从支撑材料中的同一文件生成，确保两者内容一致。}

% Example after the real file exists:
% \lstinputlisting[language=Matlab,caption={主程序}]{code/main.m}

\end{document}
